% Generated draft: requires mathematical and source review.
\documentclass[11pt]{article}
\usepackage[margin=25mm]{geometry}
\usepackage{fontspec}
\setmainfont{Latin Modern Roman}
\usepackage{amsmath,amssymb,mathtools,unicode-math}
\setmathfont{Latin Modern Math}
\usepackage{enumitem,needspace,xurl,hyperref}
\setlength{\parindent}{0pt}
\setlength{\parskip}{5pt}
\setlength{\emergencystretch}{3em}
\newcommand{\N}{\mathbb{N}}
\newcommand{\Z}{\mathbb{Z}}
\newcommand{\Q}{\mathbb{Q}}
\newcommand{\R}{\mathbb{R}}
\newcommand{\C}{\mathbb{C}}
\newcommand{\F}{\mathbb{F}}
\newcommand{\A}{\mathbb{A}}
\newcommand{\PP}{\mathbb{P}}
\newcommand{\E}{\mathbb{E}}
\newcommand{\eps}{\varepsilon}
\newcommand{\dd}{\,\mathrm{d}}
\begin{document}
\section*{UnsolvedMath ID: 30003690}
\textbf{Source problem code:} catalogue entry
\subsection{Definitions and mathematical statement}
The supplied record concatenates two independent questions. Its statement status is \texttt{unrecoverable}; the following records both questions under explicit conventional interpretations, without promoting them to a verified clean statement.

\begin{enumerate}
\item A real \(x\in[0,1)\) is Martin--L\"of random for Lebesgue measure if, for every uniformly effectively open sequence \((U_k)_{k\geq1}\) with \(\lambda(U_k)\leq 2^{-k}\), one has \(x\notin\bigcap_{k\geq1}U_k\). Here effectively open means a computably enumerable union of intervals with rational endpoints, relative to \([0,1)\).

For \(t\in\R\), let \(\{t\}=t-\lfloor t\rfloor\). Set \(y_j=\{2^j x\}\), for \(j\geq1\), and define the star discrepancy of the first \(N\) terms by
\[
D_N^*(x)=\sup_{0\leq t\leq1}\left|\frac{1}{N}\#\{1\leq j\leq N:y_j<t\}-t\right|.
\]
Under this interpretation, the first question asks whether there exist a Martin--L\"of random \(x\in[0,1)\), a constant \(C>0\), and an integer \(N_0\geq2\) such that
\[
\forall N\geq N_0,\qquad D_N^*(x)\leq C\frac{\log N}{N}.
\]
The source does not specify the domain and randomness measure, reduction modulo one, or the discrepancy convention. The definitions above make the usual unit-interval interpretation explicit. If discrepancy instead means the supremum over all half-open subintervals \([a,b)\subseteq[0,1)\), its value \(D_N\) satisfies \(D_N^*\leq D_N\leq2D_N^*\), so that choice gives the same existence question with this asymptotic bound. Other intended conventions require source clarification.

\item For this reconstruction let \(\N=\{0,1,2,\ldots\}\). For an integer \(n\geq1\) and an infinite set \(H\subseteq\N\), define
\[
\operatorname{FS}_n(H)=\{h_1+\cdots+h_n:h_1,\ldots,h_n\in H,\ h_1<\cdots<h_n\}.
\]
An instance of the exact-length, two-color restriction of Hindman's theorem is a total coloring \(c:\N\to\{0,1\}\). A solution is an infinite \(H\subseteq\N\) for which \(c\) is constant on \(\operatorname{FS}_n(H)\). An instance is computable when \(c\) is computable. A set \(H\) is \(\Sigma^0_n\) when its membership predicate is definable by an unrelativized arithmetical \(\Sigma^0_n\) formula, with \(n\) alternating blocks of number quantifiers beginning existentially and a computable matrix; unused quantifier blocks are permitted.

Under these conventions, the second question asks for which integers \(n\geq1\) the following holds:
\[
\begin{gathered}
\exists\text{ total computable }c:\N\to\{0,1\}\quad
\forall H\subseteq\N,\\
\left[
H\text{ is infinite}\ \land\
\exists i\in\{0,1\}\ \forall s\in\operatorname{FS}_n(H),\ c(s)=i
\right]
\ \Longrightarrow\ H\notin\Sigma^0_n.
\end{gathered}
\]
Here sums use exactly \(n\) distinct elements, rather than allowing repeated summands or all lengths at most \(n\). The source does not explicitly specify these conventions, whether natural numbers include zero, or the representation of a solution; these are stated interpretations requiring review.
\end{enumerate}

The supplied dated curation labels the composite record partially solved and reports a computable exact-two-sum, two-color instance with no \(\Sigma^0_2\) solution. It describes the general \(n\) question and the Martin--L\"of randomness question as unresolved. These are supplied literature assessments, not independently verified mathematical or current-status claims.
\subsection{Short English statement}
This composite record asks two independent questions: whether a Martin--L\"of random real can have discrepancy bounded by \(C\log N/N\) along its doubling orbit, and which exact sum lengths \(n\) admit a computable two-coloring with no infinite \(\Sigma^0_n\) homogeneous solution. The displayed formulations state conventional interpretations explicitly; the source marks the composite statement unrecoverable, and human mathematical review is required.
\subsection{Sources}
\begin{itemize}
\item Supplied local registry record: permanent UnsolvedMath ID \(30003690\), exact source problem code \texttt{catalogue entry}, title ``Composite of Randomness and Hindman-Theorem Questions,'' in the Oberwolfach Reports Open Problems collection. No individual proposer is supplied; the proposed year is recorded as 2018.
\item Supplied report citation: \emph{Computability Theory (2019)}, Oberwolfach Reports, DOI 10.4171/owr/2018/1, pp.\ 5--41. \url{https://doi.org/10.4171/owr/2018/1}. The background attributes the questions to workshop open-problem discussions documented in Oberwolfach Reports. The report title's year and the DOI's year component are retained as supplied.
\item Supplied literature reference: ``Csima et al., exact-two-sum Hindman theorem.'' \url{https://math.uchicago.edu/~drh/Papers/hindman.html}. The record supplies no fuller author list or bibliographic details.
\item Supplied literature reference: ``Oberwolfach report.'' \url{https://doi.org/10.4171/owr/2018/1}.
\item Literature assessment dated exactly \texttt{2026-08-21}: supplied dated curation, not independently verified. It labels the record partially solved, reports the exact-two-sum two-color result with no \(\Sigma^0_2\) solution, and describes the general \(n\) and Martin--L\"of discrepancy questions as unresolved. The background expressly requests independent verification of time-sensitive status and cited claims.
\item Supplied statement-verification note: the official open-problem section is reported to show that the extraction concatenates independent questions; no separate clean statement is promoted, and the original remains visible for provenance. This verification account is supplied metadata, not a claim that the sources were visited for this draft.
\end{itemize}
\end{document}
